\answerAnchor{MATH1-CALC-0081}
\begin{solutionBox}
\textbf{A}\quad 根号条件给
\[
D=\{(x,y):x^2+4y^2\le1\},
\]
是含边界的椭圆盘；等号处函数值为 \(0\)。

\textbf{B}\quad 对数与分母根式分别要求
\[
y-x>0,\qquad 9-x^2-y^2>0.
\]
故定义域是开圆盘 \(x^2+y^2<9\) 与半平面 \(y>x\) 的交。

\textbf{C}\quad 定义域为 \(\mathbb R^2\)。水平面 \(z=c\) 与图形相交时
\[
x^2+y^2=4-c,
\]
故 \(c\le4\) 时为半径 \(\sqrt{4-c}\) 的圆（\(c=4\) 退化为点），图形是
开口向下、顶点 \((0,0,4)\) 的旋转抛物面。

\textbf{M}\quad 圆盘 \(x^2+y^2\le a\) 上 \(x+y\) 的最大值为
\(\sqrt{2a}\)（需 \(a\ge0\)）。存在点使 \(x+y>1\) 当且仅当
\(\sqrt{2a}>1\)，所以 \(a>1/2\)。严格不等号排除了 \(a=1/2\)。
\end{solutionBox}

\answerAnchor{MATH1-CALC-0082}
\begin{solutionBox}
\textbf{A}\quad 令 \(r=\sqrt{x^2+y^2}\)，则
\[
0\le2x^2+3y^2\le3r^2\to0,
\]
故极限为 \(0\)。

\textbf{B}\quad
\[
\left|\frac{x^2y}{r}\right|
\le\frac{r^3}{r}=r^2\to0.
\]

\textbf{C}\quad \(x^4+y^4\le(x^2+y^2)^2=r^4\)，故原式绝对值不超过
\(r^3\)，极限为 \(0\)。

\textbf{M}\quad 有
\[
|x^2+3y^2|\le3(x^2+y^2)=3r^2.
\]
给定 \(\varepsilon>0\)，取
\(\delta=\sqrt{\varepsilon/3}\)。若 \(0<r<\delta\)，则
\(|x^2+3y^2|<3\delta^2=\varepsilon\)，定义证明完成。
\end{solutionBox}

\answerAnchor{MATH1-CALC-0083}
\begin{solutionBox}
\textbf{A}\quad 沿 \(y=0\) 函数值恒为 \(1\)，沿 \(y=x\) 函数值恒为
\(3/2\)，故极限不存在。

\textbf{B}\quad 分母中的两个平方项提示令
\[
y-\sin x=kx^2,\qquad\text{即 }y=\sin x+kx^2.
\]
该曲线随 \(x\to0\) 趋于原点，且代入后函数值为
\[
\frac{x^2(kx^2)}{x^4+k^2x^4}=\frac{k}{1+k^2}.
\]
取 \(k=0,1\) 得到不同值，故极限不存在。

\textbf{C}\quad 令 \(t\to0^+\)。两列点
\((t,t)\) 与 \((-t,t)\) 都在 \(D\) 内，且函数值分别为 \(1\) 与 \(-1\)。
故相对给定定义域的极限不存在。注意 \(D\) 中除原点外恒有 \(y>0\)，
所以分式在所考察的定义域内有意义。

\textbf{M}\quad 取
\[
P_n=\left(\frac1{\sqrt{\pi/2+2n\pi}},0\right),\qquad
Q_n=\left(\frac1{\sqrt{3\pi/2+2n\pi}},0\right).
\]
两列点都趋于原点，但函数值分别恒为 \(1\) 与 \(-1\)，故二重极限不存在。
\end{solutionBox}

\answerAnchor{MATH1-CALC-0084}
\begin{solutionBox}
\textbf{A}\quad
\[
\left|\frac{x^3}{x^2+y^2}\right|
\le\frac{r^3}{r^2}=r\to0.
\]

\textbf{B}\quad \(|x|,|y|\le r\)，所以分子不超过 \(r^6\)，分母为
\(r^4\)，原式绝对值不超过 \(r^2\)，极限为 \(0\)。

\textbf{C}\quad 令 \(r=\sqrt{x^2+y^2}\)，原式为
\(\sin(r^2)/r\)。由 \(|\sin(r^2)|\le r^2\)，绝对值不超过 \(r\)，
故极限为 \(0\)。

\textbf{M}\quad 极坐标下
\[
f_\alpha=r^\alpha\cos2\theta.
\]
当 \(\alpha>0\) 时 \(|f_\alpha|\le r^\alpha\to0\)；当 \(\alpha=0\) 时
沿两坐标轴分别为 \(1,-1\)；当 \(\alpha<0\) 时沿 \(x\) 轴绝对值趋于无穷。
故仅 \(\alpha>0\) 时极限存在且为 \(0\)。
\end{solutionBox}

\answerAnchor{MATH1-CALC-0085}
\begin{solutionBox}
\textbf{A}\quad 指数函数、多项式及其和在全平面连续，故 \(f\) 在
\(\mathbb R^2\) 连续。

\textbf{B}\quad 令 \(r=\sqrt{x^2+y^2}\)，则
\[
\left|\frac{xy^2}{x^2+y^2}\right|
\le\frac{r^3}{r^2}=r\to0=f(0,0),
\]
故原点连续。

\textbf{C}\quad \(f\) 是多项式，因而连续；三角形 \(D\) 由非严格线性不等式
围成，闭且有界。由连续函数在有界闭区域上的最值定理，\(f\) 在 \(D\) 上取得
最大值和最小值。

\textbf{M}\quad 删除“闭”：在
\(D=[0,1)\times[0,1]\) 上，连续函数 \(f=x+y\) 的上确界为 \(2\)，
但因 \(x<1\) 而取不到。删除“连续”：在闭圆盘上令
\[
g(x,y)=
\begin{cases}(1-x^2-y^2)^{-1},&x^2+y^2<1,\\0,&x^2+y^2=1,\end{cases}
\]
则 \(g\) 在边界不连续且无界，因而无最大值。
\end{solutionBox}

\answerAnchor{MATH1-CALC-0086}
\begin{solutionBox}
\textbf{A}\quad
\[
f_x=3x^2y-y\sin(xy),\qquad
f_y=x^3-x\sin(xy).
\]

\textbf{B}\quad 沿 \(x\) 方向 \(f(h,0)=0\)，故 \(f_x(0,0)=0\)；
沿 \(y\) 方向 \(f(0,h)=h^2\)，差商为 \(h\to0\)，故 \(f_y(0,0)=0\)。

\textbf{C}\quad 对 \(h\ne0\)，\(f(h,0)=1\)，故
\([f(h,0)-f(0,0)]/h=1/h\)，\(f_x(0,0)\) 不存在；同理
\(f(0,h)=-1\)，所以 \(f_y(0,0)\) 也不存在。沿坐标轴函数值不趋于 \(0\)，
函数在原点不连续。

\textbf{M}\quad
\[
\frac{f_\alpha(h,0)-f_\alpha(0,0)}h
=\frac{|h|^\alpha}{h}
=\operatorname{sgn}(h)|h|^{\alpha-1}.
\]
仅当 \(\alpha>1\) 时极限为 \(0\)；\(\alpha=1\) 两侧相反，
\(0<\alpha<1\) 时发散。因此 \(x\) 偏导存在当且仅当 \(\alpha>1\)。
\end{solutionBox}

\answerAnchor{MATH1-CALC-0087}
\begin{solutionBox}
\textbf{A}\quad
\[
f_{xx}=-y^2\sin(xy),\quad
f_{xy}=\cos(xy)-xy\sin(xy),
\]
\[
f_{yx}=\cos(xy)-xy\sin(xy),\quad
f_{yy}=-x^2\sin(xy).
\]

\textbf{B}\quad 由
\(f_x=2xy^3+e^xy\) 得 \(f_{xy}=6xy^2+e^x\)；
由 \(f_y=3x^2y^2+e^x\) 得 \(f_{yx}=6xy^2+e^x\)，两者一致。

\textbf{C}\quad 先沿 \(x\) 定义偏导：
\[
f_x(0,y)=\lim_{h\to0}\frac{f(h,y)}h
=\lim_{h\to0}\frac{y^3}{h^2+y^2}=y
\]
（\(y\ne0\)，而 \(y=0\) 时也为 \(0\)），故 \(f_{xy}(0,0)=1\)。
另一方面
\[
f_y(x,0)=\lim_{h\to0}\frac{f(x,h)}h
=\lim_{h\to0}\frac{xh^2}{x^2+h^2}=0,
\]
故 \(f_{yx}(0,0)=0\)。这说明无连续性条件不能交换混合偏导。

\textbf{M}\quad \(C^2\) 条件要求
\[
(f_x)_y=a+2xy=(f_y)_x=2+2xy,
\]
所以 \(a=2\)。可取
\[
f(x,y)=2xy+\frac12x^2y^2+C.
\]
\end{solutionBox}

\answerAnchor{MATH1-CALC-0088}
\begin{solutionBox}
\textbf{A}\quad
\[
f_x=2xy,\qquad f_y=x^2+3y^2.
\]
在 \((1,-1)\) 处分别为 \(-2,4\)，故 \(dz=-2\,dx+4\,dy\)。

\textbf{B}\quad 在 \((1,1)\) 线性化：
\[
f(1,1)=\sqrt7,\qquad f_x(1,1)=f_y(1,1)=-\frac1{\sqrt7}.
\]
取 \(\Delta x=0.02,\Delta y=0.01\)，估值为
\[
\sqrt7-\frac{0.03}{\sqrt7}.
\]

\textbf{C}\quad 原点两偏导为 \(0\)。又
\[
|x^2+xy+y^2|
\le x^2+\frac{x^2+y^2}{2}+y^2
=\frac32r^2,
\]
故该增量除以 \(r\) 趋于 \(0\)，函数在原点可微且 \(df=0\)。

\textbf{M}\quad
\[
dV=2\pi rh\,dr+\pi r^2\,dh.
\]
精确增量中另有
\(\pi h(dr)^2+2\pi r\,dr\,dh+\pi(dr)^2dh\)，前两项为二阶，
最后一项为三阶。
\end{solutionBox}

\answerAnchor{MATH1-CALC-0089}
\begin{solutionBox}
\textbf{A}\quad “偏导存在推出连续”为假；“偏导在邻域存在并在该点连续推出可微”
为真；“可微推出偏导存在”为真。第二句必须补足邻域与求值点条件。

\textbf{B}\quad 沿两坐标轴函数恒为 \(0\)，故原点两个偏导都存在且等于
\(0\)。但沿 \(y=x\) 有
\[
f(x,x)=\frac{x^4}{2x^4}=\frac12,
\]
不趋于 \(f(0,0)=0\)，所以函数在原点不连续。

\textbf{C}\quad 有 \(|f|\le|y|\le r\)，故连续；由坐标轴差商得
\[
f_x(0,0)=0,\qquad f_y(0,0)=1.
\]
候选线性主部是 \(y\)。又
\[
f(x,y)-y=-\frac{x^2y}{x^2+y^2}.
\]
沿 \(y=x\)，余项商为
\[
\frac{-x/2}{\sqrt2|x|},
\]
不趋于 \(0\)，所以函数不可微。

\textbf{M}\quad 令 \(r=\sqrt{x^2+y^2}\)。由
\[
|f(x,y)|\le x^2\le r^2
\]
知 \(|f|/r\le r\to0\)，故原点可微且微分为零。对 \(x\ne0\)，
\[
f_x=2x\sin(1/x)-\cos(1/x).
\]
而按定义 \(f_x(0,y)=0\)。令 \(x\to0\) 时余弦项不收敛，因此
\(f_x\) 在原点不连续。
\end{solutionBox}

\answerAnchor{MATH1-CALC-0090}
\begin{solutionBox}
\textbf{A}\quad \(f=r^{4/3}\)，原点两偏导为 \(0\)，且
\[
\frac{|f|}{r}=r^{1/3}\to0.
\]
故原点可微且全微分为零。

\textbf{B}\quad 由 \(2|xy|\le x^2+y^2=r^2\) 得
\[
|f(x,y)|\le\frac r2\to0,
\]
故函数连续；沿两坐标轴函数恒为 \(0\)，所以原点两偏导为 \(0\)。
沿 \(y=x\)，
\[
f(x,x)=\frac{|x|}{\sqrt2},\qquad
\frac{f(x,x)}{\sqrt{2x^2}}=\frac12,
\]
余项商不趋于 \(0\)，故不可微。

\textbf{C}\quad 线性主部候选为 \(3x+2y\)。余项绝对值不超过
\(r^3\)，除以 \(r\) 后不超过 \(r^2\to0\)，故函数可微，
\[
df=3\,dx+2\,dy.
\]

\textbf{M}\quad 沿 \(x\) 轴的偏导差商为
\[
\frac{|h|^{2\alpha}}h,
\]
故 \(\alpha\le1/2\) 时 \(x\) 偏导不存在，函数不可能可微。
当 \(\alpha>1/2\) 且 \(r<1\) 时，\(y^4\le y^2\)，从而
\[
0\le\frac{(x^2+y^4)^\alpha}{r}
\le r^{2\alpha-1}\to0.
\]
因此函数在原点可微当且仅当 \(\alpha>1/2\)，此时全微分为零。
\end{solutionBox}

\answerAnchor{MATH1-CALC-0091}
\begin{solutionBox}
\textbf{A}\quad 令 \(u=x-y,v=x^2+y^2\)，则
\[
z_x=f_u+2xf_v,\qquad z_y=-f_u+2yf_v,
\]
其中偏导在 \((x-y,x^2+y^2)\) 取值。

\textbf{B}\quad 令 \(u=xy,v=yz,s=zx\)。只有 \(u,v\) 依赖 \(y\)，故
\[
w_y=xf_u+zf_v.
\]

\textbf{C}\quad 令 \(u=e^x\cos y,v=e^x\sin y\)，则
\[
z_x=uf_u+vf_v,\qquad z_y=-vf_u+uf_v.
\]

\textbf{M}\quad
\[
z_x=f_u+f_v,\qquad z_y=f_u-f_v,
\]
所以
\[
z_x^2-z_y^2=(f_u+f_v)^2-(f_u-f_v)^2=4f_uf_v.
\]
\(f_u,f_v\) 均在 \((x+y,x-y)\) 取值。
\end{solutionBox}

\answerAnchor{MATH1-CALC-0092}
\begin{solutionBox}
\textbf{A}\quad 令 \(u=x^2+2y^2\)，则
\[
z_{xx}=4x^2f''(u)+2f'(u),\qquad
z_{yy}=16y^2f''(u)+4f'(u).
\]

\textbf{B}\quad 令 \(u=2x+y,v=x-2y\)。因中间变量线性，
\[
z_{xx}=4f_{uu}+4f_{uv}+f_{vv},
\]
其中使用了相应条件下 \(f_{uv}=f_{vu}\)。

\textbf{C}\quad 链式法则给
\[
z_x=ye^{xy},\qquad z_{xy}=e^{xy}+xye^{xy}.
\]
直接对 \(e^{xy}\) 求导得到同式。

\textbf{M}\quad \(u_x=2x+1,u_{xx}=2\)，故
\[
z_{xx}=f''(u)(2x+1)^2+2f'(u).
\]
所给错误式漏掉了 \(f'(u)u_{xx}\)。
\end{solutionBox}

\answerAnchor{MATH1-CALC-0093}
\begin{solutionBox}
\textbf{A}\quad
\[
du=2dx-dy,\qquad dv=dx+3dy,
\]
所以
\[
dz=(2f_u+f_v)dx+(-f_u+3f_v)dy.
\]

\textbf{B}\quad \(dz=v^2du+2uv\,dv\)，代入
\(du=dx+dy,dv=dx-dy\) 即得
\[
dz=[v^2+2uv]dx+[v^2-2uv]dy.
\]
把 \(u=x+y,v=x-y\) 回代，与直接展开求导一致。

\textbf{C}\quad \(z=r^2\cos2\theta\)，故
\[
dz=2r\cos2\theta\,dr-2r^2\sin2\theta\,d\theta.
\]

\textbf{M}\quad 直接由 \(z=x^4+y^2\) 得
\[
d^2z=12x^2(dx)^2+2(dy)^2.
\]
完整复合微分应为
\[
d^2z=2(du)^2+2u\,d^2u+d^2v.
\]
代入 \(du=2x\,dx,d^2u=2(dx)^2,d^2v=2(dy)^2\) 后恰得直接结果。
若只把外层二阶形式 \(2(du)^2\) 代入，就会漏掉
\(4x^2(dx)^2+2(dy)^2\)。
\end{solutionBox}

\answerAnchor{MATH1-CALC-0094}
\begin{solutionBox}
\textbf{A}\quad 令 \(F=x^2+2y^2+z^2-4\)。在 \((1,1,1)\)，
\(F_z=2\ne0\)，故
\[
z_x=-\frac{2x}{2z}=-1,\qquad
z_y=-\frac{4y}{2z}=-2.
\]

\textbf{B}\quad 令 \(F=xz+e^z+y^2\)，则
\[
z_x=-\frac{z}{x+e^z},\qquad
z_y=-\frac{2y}{x+e^z}.
\]

\textbf{C}\quad 全微分为
\[
(2x+yz)dx+(2y+xz)dy+xy\,dz=0.
\]
当 \(xy\ne0\) 时
\[
dz=-\frac{2x+yz}{xy}dx-\frac{2y+xz}{xy}dy.
\]

\textbf{M}\quad 令 \(F=z^3-x^2-y^2\)，则原点有 \(F_z=3z^2=0\)，
隐函数偏导公式的分母为零。实际实分支唯一：
\[
z=(x^2+y^2)^{1/3}=r^{2/3}.
\]
但沿 \(x\) 轴的差商为 \(|h|^{2/3}/h\)，不收敛，故该唯一分支在原点
不可微。这个例子说明 \(F_z=0\) 只使定理失效，实际问题既可能多分支，
也可能唯一而不可微。
\end{solutionBox}

\answerAnchor{MATH1-CALC-0095}
\begin{solutionBox}
\textbf{A}\quad 求导得
\[
u'+2v'=1,\qquad3u'-v'=2x.
\]
解得
\[
u'=\frac{1+4x}{7},\qquad v'=\frac{3-2x}{7}.
\]

\textbf{B}\quad
\[
2u\,u'+v'=1,\qquad u'+2v\,v'=e^x.
\]
关于 \((u,v)\) 的 Jacobian 为 \(4uv-1\)，非零时该系统唯一可解。

\textbf{C}\quad 对 \(x\) 求导得
\[
\begin{pmatrix}1&1\\1&-1\end{pmatrix}
\binom{u_x}{v_x}
=-\binom{y}{2x},
\]
所以
\[
\binom{u_x}{v_x}
=\binom{-x-y/2}{x-y/2}.
\]
对 \(y\) 求导同理：
\[
\begin{pmatrix}1&1\\1&-1\end{pmatrix}
\binom{u_y}{v_y}
=-\binom{x}{-1},
\qquad
\binom{u_y}{v_y}
=\binom{(1-x)/2}{-(1+x)/2}.
\]
直接解原方程得
\(u=-x^2/2-xy/2+y/2,\ v=x^2/2-xy/2-y/2\)，求偏导与上式一致。

\textbf{M}\quad 关于 \((u,v)\) 的 Jacobian 为
\[
\begin{vmatrix}3u^2&0\\0&1\end{vmatrix}=3u^2,
\]
在原点退化。方程虽有唯一实分支
\(u=\sqrt[3]x,\ v=y\)，但 \(u\) 在 \(x=0\) 的差商为
\(h^{-2/3}\)，不收敛为有限值，故不存在局部可微函数组。
\end{solutionBox}

\answerAnchor{MATH1-CALC-0096}
\begin{solutionBox}
\textbf{A}\quad 三项总次数均为 \(4\)，且
\[
f(tx,ty)=t^4f(x,y),
\]
所以是四次齐次函数。

\textbf{B}\quad \(f\) 为四次齐次函数，由 Euler 公式
\[
xf_x+yf_y=4f=4(x^2+y^2)^2.
\]

\textbf{C}\quad 对加权缩放关系关于 \(t\) 求导：
\[
a t^{a-1}x\,f_x(t^ax,t^by)
+b t^{b-1}y\,f_y(t^ax,t^by)
=m t^{m-1}f(x,y).
\]
令 \(t=1\)，得加权 Euler 公式
\[
a x f_x+b y f_y=m f.
\]

\textbf{M}\quad 从
\(f(tx,ty)=t^m f(x,y)\) 出发，固定 \(t>0\)，分别对 \(x,y\) 求导：
\[
t f_x(tx,ty)=t^m f_x(x,y),\qquad
t f_y(tx,ty)=t^m f_y(x,y).
\]
约去 \(t\) 得
\[
f_x(tx,ty)=t^{m-1}f_x(x,y),\qquad
f_y(tx,ty)=t^{m-1}f_y(x,y),
\]
所以两个一阶偏导都是 \(m-1\) 次齐次函数。
\end{solutionBox}

\answerAnchor{MATH1-CALC-0097}
\begin{solutionBox}
\textbf{A}\quad 单位方向为 \((1/\sqrt5,2/\sqrt5)\)，而
\(\nabla f(1,1)=(2,4)\)，故
\[
D_{\ell}f=\frac{2+8}{\sqrt5}=2\sqrt5.
\]
直接代入参数差商也得到同值。

\textbf{B}\quad 单位方向为
\((1/2,\sqrt3/2)\)，\(\nabla f(2,-1)=(-1,2)\)，故
\[
D_\ell f=-\frac12+\sqrt3.
\]

\textbf{C}\quad 对单位方向 \((a,b)\)，
\[
\frac{f(ta,tb)}t
=\frac{|t|\sqrt{a^2+4b^2}}t.
\]
正负 \(t\) 给出相反值，且系数非零，故按双侧定义任一方向导数都不存在。

\textbf{M}\quad 沿单位方向 \((a,b)\)，当 \(b\ne0\) 时，
\[
\frac{f(ta,tb)}t
=\frac{a^2b}{t^2a^4+b^2}\longrightarrow\frac{a^2}{b};
\]
当 \(b=0\) 时路径上函数恒为 \(0\)，方向导数也为 \(0\)。
因此每个固定方向的方向导数都存在。但沿抛物线 \(y=x^2\) 有
\[
f(x,x^2)=\frac12,
\]
函数在原点甚至不连续，因而不可微。
\end{solutionBox}

\answerAnchor{MATH1-CALC-0098}
\begin{solutionBox}
\textbf{A}\quad
\[
\nabla f=(2x+2y,2x+6y),
\]
在 \((1,-1)\) 为 \((0,-4)\)。最大方向导数为 \(4\)，方向为 \((0,-1)\)。

\textbf{B}\quad \(\nabla T(1,2)=(6,4)\)，模为 \(2\sqrt{13}\)。
最速下降方向为
\[
\left(-\frac3{\sqrt{13}},-\frac2{\sqrt{13}}\right),
\]
最小方向导数为 \(-2\sqrt{13}\)。

\textbf{C}\quad 展开或直接求偏导可得 \(\nabla f(0,0)=0\)。
但沿 \(y=0\) 有 \(f(x,0)=2x^4>0\)，而沿
\(y=\tfrac32x^2\) 有
\[
f\left(x,\tfrac32x^2\right)=-\frac14x^4<0
\qquad(x\ne0).
\]
任意小邻域内均有正、负函数值，故原点不是极值点。

\textbf{M}\quad 令 \(F=x^2+2y^2\)，则
\(\nabla F(1,1)=(2,4)\) 是法向量。切线为
\[
2(x-1)+4(y-1)=0,
\]
即 \(x+2y-3=0\)。
\end{solutionBox}

\answerAnchor{MATH1-CALC-0099}
\begin{solutionBox}
\textbf{A}\quad 令 \(F=x^2+2y^2+z^2-4\)，点满足方程且
\(\nabla F(1,1,1)=(2,4,2)\ne0\)。切平面为
\[
(x-1)+2(y-1)+(z-1)=0.
\]
法线为 \((x,y,z)=(1,1,1)+t(1,2,1)\)。

\textbf{B}\quad 令 \(F=x^2-2y^2-z\)，在点的法向量为
\((2,-4,-1)\)。切平面为
\[
2(x-1)-4(y-1)-(z+1)=0.
\]

\textbf{C}\quad 曲面法向量为 \((2x,2y,-1)\)。与
\((4,2,-1)\) 平行时由第三分量知比例系数为 \(1\)，故
\[
x=2,\qquad y=1,\qquad z=x^2+y^2=5.
\]

\textbf{M}\quad 令 \(F=x^4+y^4-z^2\)，则
\(\nabla F(0,0,0)=0\)，所以梯度公式只给恒等式 \(0=0\)，不能据此
判断切平面。实际曲面有两支
\[
z_\pm=\pm\sqrt{x^4+y^4}.
\]
令 \(r=\sqrt{x^2+y^2}\)，则
\[
|z_\pm|\le x^2+y^2=r^2=o(r),
\]
故两支在原点都可微且全微分为零，切平面同为 \(z=0\)。
因此整个曲面在一阶意义下仍有唯一切平面 \(z=0\)；梯度为零只说明公式失效，
不能单凭它断言切平面不存在。
\end{solutionBox}
